\begin{lemma}[Proved: a Bockstein parity invariant of free involutions]
\label{C:Bockstein_parity}
Let $\mathcal E$ be the exact category of finite free
$\mathbb Z/4$-modules with an involution $T$. Set
$D=T-1$, $\overline M=M/2M$, and
\[
 H(M)=\ker(\overline D:\overline M\to\overline M)\big/
                         \operatorname{im}\overline D.
\]
There is a natural endomorphism $N_M$ of this $\mathbb F_2$-space:
if $\bar v$ is a cycle, choose a lift $v\in M$ and write $Dv=2w$;
then $N_M[\bar v]=[\bar w]$.
Let $f\in\mathbb F_2[t]$ be any monic irreducible polynomial satisfying
$f(t+1)=f(t)$, and let $H(M)_f$ denote the $f$-primary part
for $N_M$. Then
\[
 \chi_f(M)=\operatorname{length}_{\mathbb F_2[t]}H(M)_f\pmod2
       \quad\in\mathbb Z/2                              \tag{C.1.2}
\]
is additive on every admissible short exact sequence. It therefore
defines a homomorphism $R_{\mathbb Z/4}(C_2)\to\mathbb Z/2$.
\end{lemma}
\begin{proof}
The identity $D^2=-2D$ implies $\overline D^2=0$.
Because $M$ is free over $\mathbb Z/4$, division by $2$ determines
$\bar w$ uniquely. Also
$D w=D^2v/2=-Dv$ modulo $2M$, so $\bar w$ is a cycle.
Replacing $v$ by $v+2u$ changes $\bar w$ by the boundary
$\overline D\bar u$. If $\bar v=\overline D\bar z$ is a boundary,
choose $v=Dz$; then $Dv/2=-Dz$ modulo $2M$, also a boundary.
This proves well-definedness, linearity and naturality of $N_M$.

Consider an admissible sequence $0\to A\to E\to B\to0$ in
$\mathcal E$. An underlying $\mathbb Z/4$-linear splitting gives
\[
 T_E=\begin{pmatrix}T_A&X\\0&T_B\end{pmatrix},\qquad
 D_AX+XD_B+2X=0.                                      \tag{C.1.3}
\]
The short exact sequence of one-periodic complexes with differentials
$\overline D$ has the exact cyclic homology sequence
\[
 H(A)\longrightarrow H(E)\longrightarrow H(B)
            \xrightarrow{\partial}H(A).
                                                               \tag{C.1.4}
\]
Here $\partial[\bar b]=[\bar X\bar b]$, and exactness follows by
the ordinary cycle-and-boundary lifting argument. The first two
maps commute with $N$. For the last map write $D_Bb=2b'$ in (C.1.3);
division by $2$ gives
\[
 N_A\partial=\partial(N_B+1).                         \tag{C.1.5}
\]
In particular the connecting map need not commute with $N$.
This twist is essential to the invariant.

Since $f(t+1)=f(t)$, every map in (C.1.4) preserves the $f$-primary
parts. The restricted cyclic sequence is still exact: decompose each
space into its primary parts, observing that $\partial$ permutes
the parts by $g(t)\mapsto g(t+1)$; only the $f$-part can map to
the $f$-part. Write $a,e,b$ for the lengths of its three spaces,
and $r,s,t$ for the lengths of the three successive images.
Exactness gives $a=t+r$, $e=r+s$, $b=s+t$.
These lengths are legitimate also at the twisted connecting map:
twisting $t$ to $t+1$ preserves $f$-primary length. Consequently
$a+e+b=2(r+s+t)$, so $e\equiv a+b\pmod2$.
This proves the asserted additivity.
\end{proof}

\begin{proposition}[Proved: flat nilpotent mixed-characteristic bases can have a generic kernel]
\label{C:flat_nilpotent_kernel}
Let
\[
 k=\mathbb Z_{(2)}[\epsilon]/(\epsilon^2),\qquad
 T=\begin{pmatrix}1&\epsilon\\\epsilon&-1\end{pmatrix},
 \qquad M=(k^2,T).
\]
Then $k$ is local, flat over $\mathbb Z_{(2)}$, and has characteristic
zero. Every integer localization $k[1/a]$ with $a\ne0$ has $K_0=\mathbb Z$,
including $k\otimes\mathbb Q$.
Nevertheless
\[
 \eta=[M]-[\mathbf1]-[\mathrm{sgn}]
       \in R_k(\mathfrak S_2)                              \tag{C.1.6}
\]
is a nonzero element of order exactly two, killed by extension to
$k\otimes\mathbb Q$. In particular $\vartheta_k$ is not surjective.
The same assertion holds with $\mathbb Z_2$ in place of
$\mathbb Z_{(2)}$.
\end{proposition}
\begin{proof}
The ring is free of rank two over $\mathbb Z_{(2)}$ and local with
maximal ideal $(2,\epsilon)$. An integer localization is either this
ring or the local ring $\mathbb Q[\epsilon]/\epsilon^2$, so its
finite projectives are free and $K_0=\mathbb Z$.
The identity $T^2=1$ follows from $\epsilon^2=0$.
There is a ring homomorphism
\[
 k\longrightarrow\mathbb Z/4,\qquad\epsilon\longmapsto2.
\]
Scalar extension is exact on the representation category because
every admissible sequence splits on the underlying projective
modules. Compose it with Lemma~\ref{C:Bockstein_parity} for
$f=t^2+t+1$.
For the resulting involution on $(\mathbb Z/4)^2$ we obtain
\[
 \overline D=0,\qquad
 N_M=\begin{pmatrix}0&1\\1&1\end{pmatrix},
       \qquad\det(t-N_M)=t^2+t+1.
\]
Thus $H(M)$ is the simple $\mathbb F_2[t]/(f)$-module and
$\chi_f(M)=1$. For the trivial and sign lines, $N$ is respectively
$0$ and $1$, and $\chi_f=0$ in both cases. This proves $\eta\ne0$.
The degree-two image of $\vartheta_k$ is generated by the trivial
line and the regular permutation module. The exact sequence
$0\to\mathbf1\to k[C_2]\to\mathrm{sgn}\to0$ identifies this
with the span of the two lines, so the same invariant proves
failure of surjectivity.

After inverting $2$, the idempotents $(1\pm T)/2$ split $M$ into
rank-one projectives. The localized ring is local, so these are
the free trivial and sign lines, and $\eta$ becomes zero.
To compute its order, the standard evaluation sequence is
\[
 0\longrightarrow M\otimes\mathrm{sgn}
 \longrightarrow k[C_2]\otimes_k k^2
 \longrightarrow M\longrightarrow0.                    \tag{C.1.7}
\]
In the middle term $C_2$ acts on its first factor. Writing its
underlying module as $k^2\oplus k^2$, the last map is
$(u,v)\mapsto u+Tv$ and the first is $v\mapsto(-Tv,v)$;
the induced action on the kernel parameter is $-T$.
The matrix
$J=\left(\begin{smallmatrix}0&1\\-1&0\end{smallmatrix}\right)$
satisfies $JT=(-T)J$ and is invertible. Hence
$M\otimes\mathrm{sgn}\simeq M$ and (C.1.7) gives
$2[M]=2[k[C_2]]=2[\mathbf1]+2[\mathrm{sgn}]$.
Together with the nonzero invariant this proves that $\eta$ has
order exactly two. The proof for $\mathbb Z_2$ is identical, using
its quotient map to $\mathbb Z/4$.
\end{proof}

\begin{corollary}[Proved: infinitely many independent flat nilpotent obstructions]
\label{C:flat_nilpotent_infinite_kernel}
For either ring $k$ in Proposition~\ref{C:flat_nilpotent_kernel}, the
kernel of
\[
 R_k(\mathfrak S_2)\longrightarrow
 R_{k\otimes\mathbb Q}(\mathfrak S_2)=\mathbb Z^2
\]
contains a subgroup isomorphic to
$\bigoplus_{r\geq0}\mathbb Z/2$.
In particular the extra degree-two representation group need not even
be finitely generated, although $k$ is noetherian, local and flat over
the mixed-characteristic DVR and every integer-localized $K_0$ is
$\mathbb Z$.
\end{corollary}
\begin{proof}
For each $r\geq0$ put $d=2^r$ and choose an element
$\alpha\in\mathbb F_{2^d}$ of trace one to $\mathbb F_2$.
Such an element exists: the trace polynomial
$x+x^2+\cdots+x^{2^{d-1}}$ is a nonzero polynomial of degree less
than $2^d$, and its values lie in $\mathbb F_2$.
The element $\alpha$ has degree exactly $d$ over $\mathbb F_2$.
Indeed if it belonged to a proper subfield $\mathbb F_{2^e}$, the
trace would be $d/e$ times its subfield trace, hence zero since $d/e$
is even. Let $q_r(u)$ be its monic irreducible polynomial.
The equation $z^2+z=\alpha$ has no solution in $\mathbb F_{2^d}$:
the map $z\mapsto z^2+z$ has kernel $\mathbb F_2$, and its image is
exactly the trace-zero subspace by dimension and the trace identity.
It follows that
\[
 f_r(t)=q_r(t^2+t)
\]
is irreducible of degree $2d$ over $\mathbb F_2$. More explicitly a
root $\beta$ generates $\alpha=\beta^2+\beta$ and a quadratic
extension of $\mathbb F_{2^d}$, so its degree is $2d$.
The polynomials $f_r$ are pairwise distinct and satisfy
$f_r(t+1)=f_r(t)$.

Lift the entries of a companion matrix of $q_r$ to $0,1$ in the
coefficient DVR, obtaining $B_r$. Its determinant is odd, since
$q_r(0)=1$, so $B_r$ is invertible. Define the free involution
\[
 T_r=\begin{pmatrix}I_d&\epsilon B_r\\
                    \epsilon I_d&-I_d\end{pmatrix},\qquad
 \eta_r=[(k^{2d},T_r)]-d[\mathbf1]-d[\mathrm{sgn}].
                                                               \tag{C.1.8}
\]
Its square is one. Under $\epsilon\mapsto2$ its Bockstein homology
is the whole $\mathbb F_2^{2d}$, with operator
\[
 N_r=\begin{pmatrix}0&\overline B_r\\I_d&I_d\end{pmatrix},
 \qquad \det(tI-N_r)=q_r(t^2+t)=f_r(t).
\]
Consequently $\chi_{f_s}(\eta_r)$ is one for $s=r$ and zero
otherwise. Each class vanishes generically by the two rational
idempotents. The invertible matrix
$J_r=\left(\begin{smallmatrix}0&B_r\\-I_d&0\end{smallmatrix}\right)$
intertwines $T_r$ and $-T_r$, so the evaluation sequence (C.1.7)
gives $2\eta_r=0$. The displayed family of parity homomorphisms
then proves that the subgroup they generate is the asserted direct
sum, without requiring a classification of all involutions.
\end{proof}

The finite certificate \texttt{code/C\_finish\_bockstein.py} constructs
the Bockstein operators from actual kernels and images over
$\mathbb F_2$. It exhausts all $260$ matrices of ranks one and two
over $\mathbb Z/4$, finding $30$ involutions, and all $1808$ block
extension candidates with endpoint ranks $(1,1),(1,2),(2,1)$.
For the $652$ admissible extensions it verifies cyclic homology
exactness, the twisted connecting identity and parity additivity.
It also verifies the rank-four evaluation sequence and the integral
dual-number identities, and the first three independent family members
of dimensions $2,4,8$, whose parity matrix is the identity; output is in
\texttt{logs/C\_finish\_bockstein.json}.
